% GENERATED FILE. Edit canonical lessons, then run npm run generate:books.
% canonical-source-hash: fnv1a64:173bd5c5f9d8a79b
% canonical-lessons: GE-C102-null, GE-C102-wort, GE-C102-antwort, GE-C102-arbeit, GE-C102-musik

\chapter{Zero, Word, Answer, Work, Music}
\label{ch:ge-zero-word-answer-work-music}
\hlchaptermodality{102}

\begin{chapteropening}
\textbf{By the end of this chapter:} \emph{I can say zero, a word, an answer, work and music.}

Complete the last lesson of chapter 102: I can say zero, a word, an answer, work and music.
\end{chapteropening}

\section[null]{null — zero}

\label{lesson:GE-C102-null}



\begin{quote}
Before the new one: say the German for a hundred, then the German for a thousand.
\end{quote}

\subsection*{You'll want to know: null}
\textbf{null} — \textquotedblleft{}zero\textquotedblright{}.

\begin{cousinweb}
\textbf{null} came in with Arabic numerals from Italian \emph{nulla}, \textquotedblleft{}nothing\textquotedblright{} $\leftarrow$ Latin \emph{nūllus}, \textquotedblleft{}none\textquotedblright{}, built from \emph{ne}, \textquotedblleft{}not\textquotedblright{}, and \emph{ūllus}, \textquotedblleft{}any\textquotedblright{} (EtymDB 2.1) — borrowed.
\end{cousinweb}

\begin{grammarlens}[title={What you've built}]
One more word for this chapter.
\end{grammarlens}

\subsection*{Your turn}
\begin{itemize}
\raggedright
  \item \emph{Say it:} \emph{null}
  \item \emph{Say it:} \emph{null}, once more
  \item \emph{Say it:} say \emph{null}
  \item \emph{Recall:} say the German for a hundred, then the German for a thousand, then say \emph{null} again
\end{itemize}

\begin{tcolorbox}[breakable,colback=teal!4,colframe=teal!35!black,title={Before you move on}]
What does null mean? (\textquotedblleft{}Zero\textquotedblright{}.) Say it once more.
\end{tcolorbox}
\section[das Wort]{das Wort — a word}

\label{lesson:GE-C102-wort}



\begin{quote}
Before the new one: say the German for a thousand, then the German for zero.
\end{quote}

\subsection*{You'll want to know: das Wort}
\textbf{das Wort} — \textquotedblleft{}a word\textquotedblright{}.

\begin{cousinweb}
\textbf{Wort} $\leftarrow$ Old High German \emph{wort} $\leftarrow$ Proto-Germanic \emph{*wurdą} — inherited, the same word as English \emph{word}. (EtymDB 2.1)
\end{cousinweb}

\begin{grammarlens}[title={What you've built}]
One more word for this chapter.
\end{grammarlens}

\subsection*{Your turn}
\begin{itemize}
\raggedright
  \item \emph{Say it:} \emph{das Wort}
  \item \emph{Say it:} \emph{das Wort}, once more
  \item \emph{Say it:} say \emph{das Wort}
  \item \emph{Recall:} say the German for a thousand, then the German for zero, then say \emph{das Wort} again
\end{itemize}

\begin{tcolorbox}[breakable,colback=teal!4,colframe=teal!35!black,title={Before you move on}]
What does das Wort mean? (\textquotedblleft{}A word\textquotedblright{}.) Say it once more.
\end{tcolorbox}
\section[die Antwort]{die Antwort — an answer}

\label{lesson:GE-C102-antwort}



\begin{quote}
Before the new one: say the German for zero, then the German for a word.
\end{quote}

\subsection*{You'll want to know: die Antwort}
\textbf{die Antwort} — \textquotedblleft{}an answer\textquotedblright{}.

\begin{cousinweb}
\textbf{Antwort} $\leftarrow$ Old High German \emph{antwurti} $\leftarrow$ Proto-Germanic \emph{*andawurdiją} (EtymDB 2.1): \emph{ant-}, \textquotedblleft{}against\textquotedblright{}, + \emph{Wort} — a word given back. English \emph{answer} pairs the same \textquotedblleft{}against\textquotedblright{} with a verb meaning \textquotedblleft{}to swear\textquotedblright{}.
\end{cousinweb}

\begin{grammarlens}[title={What you've built}]
One more word for this chapter.
\end{grammarlens}

\subsection*{Your turn}
\begin{itemize}
\raggedright
  \item \emph{Say it:} \emph{die Antwort}
  \item \emph{Say it:} \emph{die Antwort}, once more
  \item \emph{Say it:} say \emph{die Antwort}
  \item \emph{Recall:} say the German for zero, then the German for a word, then say \emph{die Antwort} again
\end{itemize}

\begin{tcolorbox}[breakable,colback=teal!4,colframe=teal!35!black,title={Before you move on}]
What does die Antwort mean? (\textquotedblleft{}An answer\textquotedblright{}.) Say it once more.
\end{tcolorbox}
\section[die Arbeit]{die Arbeit — work}

\label{lesson:GE-C102-arbeit}



\begin{quote}
Before the new one: say the German for a word, then the German for an answer.
\end{quote}

\subsection*{You'll want to know: die Arbeit}
\textbf{die Arbeit} — \textquotedblleft{}work\textquotedblright{}.

\begin{cousinweb}
\textbf{Arbeit} $\leftarrow$ Old High German \emph{arbeit} $\leftarrow$ Proto-Germanic \emph{*arbaidiz}, \textquotedblleft{}toil, hardship\textquotedblright{} — inherited. (EtymDB 2.1) Work began as hardship.
\end{cousinweb}

\begin{grammarlens}[title={What you've built}]
One more word for this chapter.
\end{grammarlens}

\subsection*{Your turn}
\begin{itemize}
\raggedright
  \item \emph{Say it:} \emph{die Arbeit}
  \item \emph{Say it:} \emph{die Arbeit}, once more
  \item \emph{Say it:} say \emph{die Arbeit}
  \item \emph{Recall:} say the German for a word, then the German for an answer, then say \emph{die Arbeit} again
\end{itemize}

\begin{tcolorbox}[breakable,colback=teal!4,colframe=teal!35!black,title={Before you move on}]
What does die Arbeit mean? (\textquotedblleft{}Work\textquotedblright{}.) Say it once more.
\end{tcolorbox}
\section[die Musik]{die Musik — music}

\label{lesson:GE-C102-musik}



\begin{quote}
Before the new one: say the German for an answer, then the German for work.
\end{quote}

\subsection*{You'll want to know: die Musik}
\textbf{die Musik} — \textquotedblleft{}music\textquotedblright{}.

\begin{cousinweb}
\textbf{Musik} $\leftarrow$ Latin \emph{mūsica} $\leftarrow$ Greek \emph{mousik\'{\={e}}}, \textquotedblleft{}the art of the Muses\textquotedblright{} (EtymDB 2.1) — borrowed, like English \emph{music}.
\end{cousinweb}

\begin{grammarlens}[title={What you've built}]
One more word for this chapter.
\end{grammarlens}

\subsection*{Your turn}
\begin{itemize}
\raggedright
  \item \emph{Say it:} \emph{die Musik}
  \item \emph{Say it:} \emph{die Musik}, once more
  \item \emph{Say it:} say \emph{die Musik}
  \item \emph{Recall:} say the German for an answer, then the German for work, then say \emph{die Musik} again
\end{itemize}

\begin{tcolorbox}[breakable,colback=teal!4,colframe=teal!35!black,title={Before you move on}]
What does die Musik mean? (\textquotedblleft{}Music\textquotedblright{}.) Say it once more.
\end{tcolorbox}
